\documentclass[10pt,letterpaper,twocolumn]{article}
\usepackage[margin=0.78in,top=0.82in,bottom=0.78in,columnsep=0.27in]{geometry}
\usepackage[T1]{fontenc}
\usepackage[utf8]{inputenc}
\usepackage{mathptmx}
\usepackage{amsmath,amssymb,booktabs,graphicx,microtype,caption,xcolor,placeins,array,multirow}
\usepackage[hidelinks]{hyperref}
\usepackage{fancyhdr}
\usepackage{enumitem}
\usepackage{float}
\usepackage{cuted}
\usepackage{needspace}
\usepackage{etoolbox}
\graphicspath{{figures/}}
\captionsetup{font=footnotesize,labelfont=bf,justification=raggedright,singlelinecheck=false,skip=3pt}
\setlength{\parindent}{1em}
\setlength{\parskip}{0pt}
\setlength{\textfloatsep}{9pt plus 2pt minus 2pt}
\setlength{\floatsep}{8pt plus 2pt minus 2pt}
\setlength{\intextsep}{7pt plus 2pt minus 2pt}
\setcounter{topnumber}{5}
\setcounter{totalnumber}{8}
\renewcommand{\topfraction}{.92}
\renewcommand{\textfraction}{.06}
\renewcommand{\floatpagefraction}{.76}
\pagestyle{fancy}
\fancyhf{}
\fancyfoot[C]{\small\thepage}
\renewcommand{\headrulewidth}{0pt}
% Strictly in-column, nonfloating figure panels matching the sample paper.
% No strip, figure*, or float; figure and caption stay together.
\newcommand{\fixedfig}[2]{%
  \par\medskip
  \Needspace{0.32\textheight}%
  \noindent\begin{minipage}{\columnwidth}
     \centering
     \includegraphics[width=\columnwidth,height=0.245\textheight,keepaspectratio]{#1}\par\vspace{3pt}
     \captionsetup{type=figure}\captionof{figure}{#2}
  \end{minipage}\par\medskip
}
\title{\bfseries\Large Numerical Artifacts and the Reliability of Chaos Diagnostics in Gravitational $N$-Body Dynamics}
\author{Fatima\\\small Independent computational astrophysics research}
\date{}
\begin{document}
\twocolumn[\begin{@twocolumnfalse}
\maketitle
\vspace{-1.4em}
\begin{abstract}
\noindent Numerical integration errors can change orbital outcomes and distort diagnostics commonly used to infer chaotic behavior. We investigate forward Euler, kick--drift--kick leapfrog, fourth-order Runge--Kutta (RK4), and adaptive methods in normalized Newtonian two-body and circular restricted three-body systems. An analytically regular control, close-encounter experiments, variational-equation calculations, and a structured ensemble of 18 initial conditions were examined. Of 108 fixed-step ensemble integrations, 102 admitted independently refined adaptive-reference comparisons. Under exploratory criteria of at most 10\% relative final-state disagreement and 5\% relative finite-time Lyapunov exponent (FTLE) disagreement, RK4 passed both criteria in 18/18 verified regular-neighborhood comparisons and 7/16 verified high-sensitivity comparisons; Euler passed neither in any verified comparison. Some runs reproduced an FTLE while disagreeing substantially in physical state. A strict relative Jacobi-drift screen below $10^{-6}$ passed 15 comparisons, all of which also passed both numerical diagnostics, while rejecting 19 otherwise acceptable comparisons. Independent DOP853--Radau comparisons supported selected finite-time calculations and exposed the limitations of invariant-only validation. The resulting reliability atlas characterizes numerical agreement, not definitive asymptotic chaos labels.
\end{abstract}
\vspace{.4em}
\noindent\textbf{Keywords:} numerical gravitational dynamics; $N$-body problem; circular restricted three-body problem; Lyapunov exponent; Jacobi invariant; convergence; numerical artifacts.
\vspace{1.0em}
\end{@twocolumnfalse}]

\section{Introduction}
Numerical solutions of gravitational equations are central to celestial mechanics, yet deterministic physical laws do not guarantee accurate numerical predictions. Discretization and roundoff errors can shift orbital phases, alter the geometry of a close encounter, and change the apparent stability of a computed trajectory. In dynamically sensitive systems, very small errors may be amplified; in regular systems, poor discretization may create unphysical secular energy changes. Distinguishing physical instability from numerical artifacts is therefore essential \cite{hairer2006geometric,wisdom1991symplectic,lichtenberg1992regular}.

Integration methods preserve different properties. Symplectic schemes preserve discrete Hamiltonian structure but are not automatically the most accurate for every observable or at every timestep. High-order adaptive schemes can accurately reproduce finite-time trajectories, but neither small local error estimates nor a conserved scalar invariant prove that a computed chaos indicator is correct \cite{hairer2006geometric,rein2015ias15}.

We ask whether apparently trustworthy gravitational simulations can fail a different numerical reliability check. The study combines analytical controls, timestep-convergence experiments, perturbed orbit measurements, close encounters, an ensemble reliability atlas, and cross-algorithm validation. The central claim is deliberately about \emph{finite-time numerical reliability}. We do not equate any positive finite-time exponent with a demonstrated positive asymptotic Lyapunov exponent.

\section{Mathematical framework}
\subsection{Newtonian equations and initial conditions}
For $N$ particles, the acceleration of body $i$ is
\begin{equation}
 \ddot{\boldsymbol r}_i=G\sum_{j\ne i}m_j\frac{\boldsymbol r_j-\boldsymbol r_i}{\lVert\boldsymbol r_j-\boldsymbol r_i\rVert^3}.
\end{equation}
We use normalized units, $G=1$. The analytically regular control is a circular two-body orbit. In the circular restricted three-body cases, a star of mass unity and a planet of mass $10^{-3}$ or $5\times10^{-3}$ follow circular barycentric motion with initial separation one. A massless test particle feels both gravitational fields without back-reacting on the primaries. Its initial radius and orbital phase parameterize the experiment.

\subsection{Tangent dynamics and FTLE}
Linearizing the gravitational acceleration yields
\begin{align}
 \delta\ddot{\boldsymbol r}_i=G\sum_{j\ne i} m_j\bigg[&\frac{\delta\boldsymbol R_{ij}}{R_{ij}^3}\\[-.5em]
 &-\frac{3\boldsymbol R_{ij}(\boldsymbol R_{ij}\!\cdot\!\delta\boldsymbol R_{ij})}{R_{ij}^5}\bigg],\nonumber
\end{align}
where $\boldsymbol R_{ij}=\boldsymbol r_j-\boldsymbol r_i$. A tangent vector is periodically normalized, and its logarithmic amplification $g_k$ is accumulated:
\begin{equation}
 \lambda_T=\frac{1}{T}\sum_{k}\log g_k.
\end{equation}
Our FTLE norm combines position and normalized velocity perturbations. For a restricted three-body test, the primaries are held unperturbed in the tangent system when explicitly isolating the massless particle's sensitivity. The FTLE depends on observation horizon, norm and initial tangent direction; its positivity at finite $T$ is not a proof of long-time chaos \cite{benettin1980lyapunov1,benettin1980lyapunov2}.

\subsection{Conserved quantities}
The circular restricted three-body equations conserve the Jacobi constant in a corotating barycentric frame:
\begin{equation}
 C_J=2\left(\frac{Gm_1}{r_1}+\frac{Gm_2}{r_2}+\frac{\omega^2(x^2+y^2)}{2}\right)-v_{\rm rot}^2.
\end{equation}
We measure sampled maximum relative drift $\epsilon_J=\max_t|C_J(t)-C_J(0)|/|C_J(0)|$. This scalar constrains a trajectory to an invariant surface but does not uniquely specify its location on that surface.

\section{Numerical implementation}
\subsection{Four integration approaches}
Forward Euler is used as a first-order baseline. Kick--drift--kick leapfrog is second-order and symplectic; its tangent equations are derived by differentiating the same discrete kick and drift operations. Explicit RK4 provides fourth-order accuracy in smooth, fixed-horizon problems. Adaptive DOP853 supplies the primary high-order reference; implicit Radau provides an algorithmically distinct check \cite{hairer2006geometric,virtanen2020scipy}.

\subsection{Controlled numerical errors}
We distinguish phase-space agreement from FTLE agreement. With $X$ denoting the star-relative position and velocity vector, we use
\begin{equation}
 \epsilon_X=\frac{\lVert X_{\rm fixed}-X_{\rm ref}\rVert}{\lVert X_{\rm ref}\rVert},\quad
 \epsilon_\lambda=\frac{|\lambda_{\rm fixed}-\lambda_{\rm ref}|}{|\lambda_{\rm ref}|}.
\end{equation}
For the ensemble, an exploratory ``both consistent'' outcome requires $\epsilon_X\le 0.10$ and $\epsilon_\lambda\le 0.05$. When the reference fails verification, the comparison is \emph{unresolved}; failures are never counted as successes.

\section{Baseline experiments and encounter resolution}
\subsection{Euler drift and orbital geometry}
A forward Euler simulation accumulated energy error monotonically, whereas reducing its timestep systematically decreased this drift (Fig.~\ref{fig:euler}). The same nominal orbital initial conditions evolved into increasingly distorted trajectories under Euler, while leapfrog preserved much more nearly closed orbits (Fig.~\ref{fig:orbits}). Angular-momentum monitoring showed that leapfrog drift was close to floating-point scales in this simple benchmark, in sharp contrast with Euler (see the angular-momentum comparison).

\fixedfig{euler_timestep_panel1.png}{\label{fig:euler}Effect of timestep on forward Euler energy error. The monotonic drift decreases as the timestep is refined, illustrating the method's systematic bias in orbital integration. (panel 1 of 2).}\fixedfig{euler_timestep_panel2.png}{Effect of timestep on forward Euler energy error. The monotonic drift decreases as the timestep is refined, illustrating the method's systematic bias in orbital integration. (panel 2 of 2).}
\fixedfig{orbit_comparison_panel1.png}{\label{fig:orbits}Star--planet orbital geometry under forward Euler (left) and leapfrog (right). Euler trajectories spiral; leapfrog orbits remain approximately closed over this comparison. (panel 1 of 2).}\fixedfig{orbit_comparison_panel2.png}{Star--planet orbital geometry under forward Euler (left) and leapfrog (right). Euler trajectories spiral; leapfrog orbits remain approximately closed over this comparison. (panel 2 of 2).}
\fixedfig{angular_momentum.png}{\label{fig:ang}Angular-momentum error of Euler, leapfrog and RK4. A logarithmic ordinate exposes orders-of-magnitude differences obscured on linear axes.}

\subsection{Close encounters and local convergence}
Close passages are challenging because derivatives of the acceleration grow rapidly as separation decreases. In the encounter experiment, a candidate test particle approached a massive planet closely enough that fixed timesteps exceeding the characteristic encounter timescale led to grossly different scattering outcomes. An adaptive reference resolved the encounter with rapidly shrinking accepted steps, and high-resolution comparisons converged on the closest-approach time and distance.

\fixedfig{adaptive_encounter_panel1.png}{\label{fig:encounter}Adaptive reference near the closest planetary encounter. The minimum is sharply localized in time, requiring integration steps much smaller than those adequate for a regular orbit. (panel 1 of 2).}\fixedfig{adaptive_encounter_panel2.png}{Adaptive reference near the closest planetary encounter. The minimum is sharply localized in time, requiring integration steps much smaller than those adequate for a regular orbit. (panel 2 of 2).}
\fixedfig{adaptive_step_panel1.png}{\label{fig:step}Adaptive timestep behavior across the encounter and in a magnified interval around the minimum separation. Accepted steps contract sharply near the most demanding part of the trajectory. (panel 1 of 2).}\fixedfig{adaptive_step_panel2.png}{Adaptive timestep behavior across the encounter and in a magnified interval around the minimum separation. Accepted steps contract sharply near the most demanding part of the trajectory. (panel 2 of 2).}
\fixedfig{rk4_failure.png}{\label{fig:fail}Fixed RK4 timestep compared with the encounter timescale: coarse steps can exceed the local resolution requirement by several orders of magnitude.}
\fixedfig{local_convergence_panel1.png}{\label{fig:order}Observed local RK4 convergence of position and velocity errors in a resolved encounter window. Approximately fifth-order \emph{local} error behavior is compatible with fourth-order global RK4 convergence. (panel 1 of 2).}\fixedfig{local_convergence_panel2.png}{Observed local RK4 convergence of position and velocity errors in a resolved encounter window. Approximately fifth-order \emph{local} error behavior is compatible with fourth-order global RK4 convergence. (panel 2 of 2).}

\subsection{Spurious orbital escape in a regular control}
A separate two-body benchmark established that coarse RK4 discretization can generate unphysical energy gains and artificial unbinding despite the analytically bound initial orbit. The apparent escape time depends strongly on timestep, emphasizing that qualitative orbital classification can be corrupted by the integrator even without physical chaos. Figure~\ref{fig:escape} illustrates the numerical transition.
\fixedfig{false_escape_panel1.png}{\label{fig:escape}Numerical orbital unbinding versus RK4 timestep. Artificial escape occurs at coarse steps in an otherwise regular, analytically bound two-body control. (panel 1 of 2).}\fixedfig{false_escape_panel2.png}{Numerical orbital unbinding versus RK4 timestep. Artificial escape occurs at coarse steps in an otherwise regular, analytically bound two-body control. (panel 2 of 2).}

\section{Finite-time sensitivity experiments}
\subsection{Regular controls and high-sensitivity candidates}
The first regular two-body FTLE decayed with integration time; its positive early value was a transient finite-horizon measurement rather than a chaos classification. Scans of restricted three-body initial radii and perturbing-planet masses revealed heterogeneous sensitivity, not a monotonic relationship with perturbation strength. A candidate at test radius 1.34 and planetary mass 0.005 exhibited stronger finite-time tangent growth than nearby cases. Another initial condition near radius 1.59 was less sensitive over the tested interval.

\fixedfig{ftle_candidate.png}{\label{fig:ftle}Adaptive FTLE histories for two three-body candidates over $T=240$. The more sensitive trajectory retains a substantially larger finite-time estimate, but both require long-time qualification.}
\fixedfig{long_growth.png}{\label{fig:growth}Accumulated tangent-vector logarithmic growth for the high-sensitivity candidate. Its approximately linear portions represent sustained \emph{finite-window} amplification, not proof of an asymptotic positive exponent.}

\subsection{Numerical predictability and convergence}
Independent adaptive trajectories from identical initial states separate over long time intervals. Their relative position difference grew approximately exponentially over a selected intermediate window, with a fitted numerical discrepancy-growth rate of about 0.0446 in normalized units. This is not identical to a maximal Lyapunov exponent, because numerical errors are injected repeatedly by the algorithms. The experiment instead illustrates a finite numerical predictability horizon.

\section{Ensemble reliability atlas}
\subsection{Design and verified denominator}
Two structured neighborhoods were sampled: nine initial conditions near radius 1.40, planet mass 0.001, and nine near radius 1.34, mass 0.005. Radius offsets were $-0.01,0,+0.01$ and phase offsets were $-5^\circ,0,+5^\circ$. Each case was integrated with Euler, leapfrog and RK4 at $\Delta t=0.04$ and $0.10$, yielding 108 attempts at $T=120$. Two refined DOP853 references were checked before assessing each fixed-step result: 102 comparisons were retained, six linked to one unverified initial configuration were unresolved.

\begin{table}[H]\centering\caption{Both-diagnostic agreement rates under the exploratory 10\% final-state and 5\% relative-FTLE limits. Fractions give passing cases over verified comparisons; missing cases are not counted as passes.}\label{tab:atlas}\small\begin{tabular}{llccc}\toprule Family & Step & Euler & Leapfrog & RK4\\\midrule Regular-looking & 0.04 & 0/9 & 9/9 & 9/9\\ Regular-looking & 0.10 & 0/9 & 0/9 & 9/9\\ High-sensitivity & 0.04 & 0/8 & 0/8 & 6/8\\ High-sensitivity & 0.10 & 0/8 & 0/8 & 1/8\\\bottomrule\end{tabular}\end{table}

\fixedfig{ensemble_ftle_panel1.png}{\label{fig:ensemble}Median relative FTLE disagreement across the ensemble. RK4 is highly accurate at the smaller timestep in both neighborhoods; increasing timestep strongly affects the high-sensitivity cases. (panel 1 of 2).}\fixedfig{ensemble_ftle_panel2.png}{Median relative FTLE disagreement across the ensemble. RK4 is highly accurate at the smaller timestep in both neighborhoods; increasing timestep strongly affects the high-sensitivity cases. (panel 2 of 2).}
\fixedfig{ensemble_scatter_panel1.png}{\label{fig:scatter}State error versus FTLE error for the two initial-condition neighborhoods. An apparently accurate FTLE can coexist with substantial final-state disagreement. (panel 1 of 2).}\fixedfig{ensemble_scatter_panel2.png}{State error versus FTLE error for the two initial-condition neighborhoods. An apparently accurate FTLE can coexist with substantial final-state disagreement. (panel 2 of 2).}
\fixedfig{threshold_heatmap.png}{\label{fig:thresh}Frequency of misleading FTLE agreement across exploratory state-error and FTLE-error thresholds. This robustness check prevents reliance on a single arbitrary cutoff.}

\subsection{Numerical outcomes}
Table~\ref{tab:atlas} and Fig.~\ref{fig:atlasplot} show a strong dependence on numerical method and dynamical regime. RK4 met both criteria in all 18 verified regular-neighborhood comparisons, but only seven of 16 verified high-sensitivity comparisons. Leapfrog passed nine of 18 regular-neighborhood comparisons at the finer step and none of the high-sensitivity comparisons; Euler passed no verified comparisons under the joint rule. These outcomes describe this sampled grid and time horizon, not universal method rankings.

\fixedfig{atlas_panel1.png}{\label{fig:atlasplot}Final numerical reliability atlas, with verified numerators/denominators by neighborhood, integration method and timestep. Darker cells indicate larger fractions in each diagnostic category; the middle panel marks FTLE agreement that conceals state error. (panel 1 of 2).}\fixedfig{atlas_panel2.png}{Final numerical reliability atlas, with verified numerators/denominators by neighborhood, integration method and timestep. Darker cells indicate larger fractions in each diagnostic category; the middle panel marks FTLE agreement that conceals state error. (panel 2 of 2).}
\fixedfig{outcome_panel1.png}{\label{fig:outcomes}Outcome distribution counting all attempted simulations, including unresolved cases. The four verified categories distinguish joint agreement, one-diagnostic disagreement, and failure of both diagnostics. (panel 1 of 2).}\fixedfig{outcome_panel2.png}{Outcome distribution counting all attempted simulations, including unresolved cases. The four verified categories distinguish joint agreement, one-diagnostic disagreement, and failure of both diagnostics. (panel 2 of 2).}
\fixedfig{robustness_panel1.png}{\label{fig:robust}Sensitivity of the joint-agreement rate to alternative state and FTLE error limits. Method rankings remain stable within the tested threshold grid, although rates are operational rather than universal. (panel 1 of 2).}\fixedfig{robustness_panel2.png}{Sensitivity of the joint-agreement rate to alternative state and FTLE error limits. Method rankings remain stable within the tested threshold grid, although rates are operational rather than universal. (panel 2 of 2).}

\section{Conservation law as a diagnostic}
\subsection{Screening behavior}
Across 102 verified comparisons, a relative Jacobi-drift threshold of $10^{-6}$ accepted 15 runs, all of which also met the state and FTLE accuracy criteria. The screen rejected 19 other runs that passed both criteria and 68 that failed at least one. Thus strict Jacobi conservation was a useful but selective positive screen \emph{within this sample}. The absence of false reassurance among only 15 accepted cases does not establish a general guarantee.

\fixedfig{conservation_screen.png}{\label{fig:jacobitab}Conservation-screen contingency table for 102 verified comparisons. Fifteen conservation-passing runs pass both numerical diagnostics, but 19 acceptable runs fail the strict conservation threshold.}
\fixedfig{conservation_scatter_panel1.png}{\label{fig:jacobiscatter}Jacobi drift against FTLE error in both neighborhoods. Error growth is often related, but neither quantity alone fully specifies the other. (panel 1 of 2).}\fixedfig{conservation_scatter_panel2.png}{Jacobi drift against FTLE error in both neighborhoods. Error growth is often related, but neither quantity alone fully specifies the other. (panel 2 of 2).}
\fixedfig{jacobi_heatmap_panel1.png}{\label{fig:jacobithreshold}Sensitivity of apparent conservation-based reassurance to alternative Jacobi and FTLE error thresholds. Percentages use all verified comparisons as the denominator. (panel 1 of 2).}\fixedfig{jacobi_heatmap_panel2.png}{Sensitivity of apparent conservation-based reassurance to alternative Jacobi and FTLE error thresholds. Percentages use all verified comparisons as the denominator. (panel 2 of 2).}

\subsection{Targeted case C12}
The numerically demanding C12 initial condition failed the original reference-verification screen despite extremely small Jacobi drift. A three-level DOP853 refinement showed that the refined--ultra-refined trajectory disagreement was materially smaller than the standard--refined disagreement, while the Jacobi drift remained near floating-point scales. This demonstrates why invariant conservation, trajectory convergence and FTLE reproducibility are complementary tests.

\fixedfig{c12_diagnostics.png}{\label{fig:c12}C12 orbital geometry, successive refinement disagreements, Jacobi drift and proximity to the perturbing planet. Narrow encounter passages coincide with rapid amplification of numerical discrepancies.}
\fixedfig{c12_ftle_panel1.png}{\label{fig:c12ftle}C12 FTLE histories and numerical disagreement under increasing adaptive accuracy. Visually overlapping sensitivity curves conceal measurable differences visible on a logarithmic scale. (panel 1 of 2).}\fixedfig{c12_ftle_panel2.png}{C12 FTLE histories and numerical disagreement under increasing adaptive accuracy. Visually overlapping sensitivity curves conceal measurable differences visible on a logarithmic scale. (panel 2 of 2).}

\section{Independent algorithmic validation}
Three representative configurations (C05, C12 and C14) were independently integrated with explicit DOP853 and implicit Radau, at standard and refined tolerances. The refined trajectories and FTLE histories agreed closely, with C12 showing the largest discrepancies. Independent algorithmic agreement strengthens finite-time validation but does not establish the asymptotic dynamical classification of the orbit.

\fixedfig{cross_solver_orbits_panel1.png}{\label{fig:crossorbit}Cross-method star-relative trajectories for C05, C12 and C14. Refined DOP853 and Radau overlap closely at the scale of the plotted orbits. (panel 1 of 3).}\fixedfig{cross_solver_orbits_panel2.png}{Cross-method star-relative trajectories for C05, C12 and C14. Refined DOP853 and Radau overlap closely at the scale of the plotted orbits. (panel 2 of 3).}\fixedfig{cross_solver_orbits_panel3.png}{Cross-method star-relative trajectories for C05, C12 and C14. Refined DOP853 and Radau overlap closely at the scale of the plotted orbits. (panel 3 of 3).}
\fixedfig{cross_solver_position_panel1.png}{\label{fig:crosspos}Absolute position disagreements for solver refinement within DOP853 and Radau and the refined cross-algorithm comparison. C12 exhibits the strongest numerical sensitivity. (panel 1 of 3).}\fixedfig{cross_solver_position_panel2.png}{Absolute position disagreements for solver refinement within DOP853 and Radau and the refined cross-algorithm comparison. C12 exhibits the strongest numerical sensitivity. (panel 2 of 3).}\fixedfig{cross_solver_position_panel3.png}{Absolute position disagreements for solver refinement within DOP853 and Radau and the refined cross-algorithm comparison. C12 exhibits the strongest numerical sensitivity. (panel 3 of 3).}
\fixedfig{cross_solver_ftle_panel1.png}{\label{fig:crossftle}Finite-time Lyapunov histories from two independent integration algorithms. Agreement is assessed quantitatively in Fig.~\ref{fig:crossdiff}, not merely from overlapping curves. (panel 1 of 3).}\fixedfig{cross_solver_ftle_panel2.png}{Finite-time Lyapunov histories from two independent integration algorithms. Agreement is assessed quantitatively in Fig.~\ref{fig:crossdiff}, not merely from overlapping curves. (panel 2 of 3).}\fixedfig{cross_solver_ftle_panel3.png}{Finite-time Lyapunov histories from two independent integration algorithms. Agreement is assessed quantitatively in Fig.~\ref{fig:crossdiff}, not merely from overlapping curves. (panel 3 of 3).}
\fixedfig{cross_solver_error_panel1.png}{\label{fig:crossdiff}Logarithmic FTLE disagreement between standard and refined solutions and between refined DOP853 and Radau. (panel 1 of 3).}\fixedfig{cross_solver_error_panel2.png}{Logarithmic FTLE disagreement between standard and refined solutions and between refined DOP853 and Radau. (panel 2 of 3).}\fixedfig{cross_solver_error_panel3.png}{Logarithmic FTLE disagreement between standard and refined solutions and between refined DOP853 and Radau. (panel 3 of 3).}
\fixedfig{cross_solver_jacobi_panel1.png}{\label{fig:crossjacobi}Sampled Jacobi drift under DOP853 and Radau. Both algorithms preserve the invariant closely over the selected finite horizon, but the level of drift differs by case and solver settings. (panel 1 of 3).}\fixedfig{cross_solver_jacobi_panel2.png}{Sampled Jacobi drift under DOP853 and Radau. Both algorithms preserve the invariant closely over the selected finite horizon, but the level of drift differs by case and solver settings. (panel 2 of 3).}\fixedfig{cross_solver_jacobi_panel3.png}{Sampled Jacobi drift under DOP853 and Radau. Both algorithms preserve the invariant closely over the selected finite horizon, but the level of drift differs by case and solver settings. (panel 3 of 3).}

\section{Discussion}
Three interlocking conclusions emerge. First, under-resolved gravitational encounters and coarse timesteps can change qualitative orbital predictions, including artificial escape in an analytically regular system. Second, numerical agreement of a scalar sensitivity indicator is not equivalent to agreement of the phase-space state, nor does close conservation of one invariant uniquely determine the trajectory. Third, numerical reliability depends on initial configuration, timestep, method and observation horizon. These findings support a multi-diagnostic verification strategy rather than any singular accuracy test.

The ensemble percentages have important limits. The initial conditions form a structured grid, not an independent random population sample. The two three-body families are designated regular-looking and high-sensitivity based on finite-time evidence; neither family is guaranteed to contain exclusively regular or chaotic solutions. The 10\% state and 5\% FTLE thresholds were chosen operationally and are not universal physical boundaries. In potentially chaotic motion, pointwise trajectories naturally lose reproducibility after finite horizons even when selected long-time statistics may remain meaningful.

\section{Conclusions}
We constructed a numerical reliability atlas connecting gravitational state error, finite-time Lyapunov disagreement and Jacobi conservation. Euler, leapfrog and RK4 differed markedly in accuracy across two restricted three-body neighborhoods, with RK4 performing best under the joint finite-horizon criteria used here. Under-resolved integrations caused misleading orbital outcomes and both amplification and suppression of chaos-related measurements. A strict invariant screen detected many inaccurate simulations but also excluded accurate ones, while C12 showed why conservation alone cannot certify pointwise convergence. Independent DOP853--Radau calculations strengthened selected reference results. Future work should investigate larger ensembles, longer horizons, multiple tangent directions and independent high-precision reference algorithms before making asymptotic regular-versus-chaotic classifications.

\section*{Data availability and declarations}
The results originate from the computational notebook and the project figures. Final publication should archive the notebook, package versions, initial conditions and objective-specific CSV exports (especially Objectives 14--18), and rerun aggregate results directly from those archived tables. This draft does not claim a public repository has already been released. AI assistance was used in code explanation, debugging and manuscript preparation; all scientific interpretations and numerical results require author verification before submission.

\begin{thebibliography}{9}
\small

\bibitem{hairer2006geometric}
E.~Hairer, C.~Lubich, and G.~Wanner,
\emph{Geometric Numerical Integration: Structure-Preserving Algorithms for Ordinary Differential Equations},
2nd ed., Springer, 2006.

\bibitem{wisdom1991symplectic}
J.~Wisdom and M.~Holman,
``Symplectic maps for the N-body problem,''
\emph{The Astronomical Journal},
vol.~102, pp.~1528--1538, 1991.

\bibitem{lichtenberg1992regular}
A.~J.~Lichtenberg and M.~A.~Lieberman,
\emph{Regular and Chaotic Dynamics},
2nd ed., Springer-Verlag, 1992.

\bibitem{rein2015ias15}
H.~Rein and D.~S.~Spiegel,
``IAS15: A fast, adaptive, high-order integrator for gravitational dynamics, accurate to machine precision over a billion orbits,''
\emph{Monthly Notices of the Royal Astronomical Society},
vol.~446, no.~2, pp.~1424--1437, 2015.

\bibitem{benettin1980lyapunov1}
G.~Benettin, L.~Galgani, A.~Giorgilli, and J.-M.~Strelcyn,
``Lyapunov Characteristic Exponents for Smooth Dynamical Systems and for Hamiltonian Systems; A Method for Computing All of Them. Part 1: Theory,''
\emph{Meccanica},
vol.~15, no.~1, pp.~9--20, 1980.

\bibitem{benettin1980lyapunov2}
G.~Benettin, L.~Galgani, A.~Giorgilli, and J.-M.~Strelcyn,
``Lyapunov Characteristic Exponents for Smooth Dynamical Systems and for Hamiltonian Systems; A Method for Computing All of Them. Part 2: Numerical Application,''
\emph{Meccanica},
vol.~15, no.~1, pp.~21--30, 1980.

\bibitem{virtanen2020scipy}
P.~Virtanen et al. (SciPy 1.0 Contributors),
``SciPy 1.0: Fundamental algorithms for scientific computing in Python,''
\emph{Nature Methods},
vol.~17, pp.~261--272, 2020.

\end{thebibliography}
\end{document}
